\subsection{Curvature and TNB Vectors}

\subsubsection{Curvature}


\content{
\begin{definition} A parametrization of a curve $\vec{r}(t)$ is called smooth if
$\vec{r}'(t)$ is continuous and $\vec{r}'(t)\neq \vec{0}$. A curve is called
smooth if it has a smooth parametrization. 
\end{definition}
}{% presentation
}{% nopresentation
}{% comments
A smooth curve has no sharp corners or cusps. When it turns the tangent vectors
turns continuously. 
}

\content{
\vspace{2in}
\begin{definition}
The curvature of a curve $\vec{r}'(t)$ is defined to be the number 
$$ \kappa = \left| \frac{d\vec{T}}{ds}\right| $$
\end{definition}
}{% presentation
}{% nopresentation
}{% comments
Draw a curve, then notice that where the curve does not turn a lot, the change
in tangent vector is very small. 
Notice that we can also use the chain rule to rewrite this in terms of the
parametrization $\vec{r}'(t)$: 
$$ \frac{d\vec{T}}{dt} = \frac{d\vec{T}}{ds}\frac{ds}{dt} $$
and we can use the fundamental theorem of calculus to find $\frac{ds}{dt}$.
}


\content{
\begin{example}
Find the curvature of a circle of radius $a$. 
\end{example}
}{% presentation
\vspace{3in}
}{% nopresentation
\vspace{2in}
}{% comments
curvature is $1/a$.
}

\subsubsection{The Tangent, Normal, and Binormal Vectors}

\content{
\begin{definition} Given a vector valued function $\vec{r}(t)$ the unit tangent
vector is given by 
$$ \vec{T}(t) = \frac{\vec{r}'(t)}{|\vec{r}'(t)|}. $$
\end{definition}
}{% presentation
\vspace{1in}
}{% nopresentation
\vspace{1in}
}{% comments
This is simply a unit vector in the same direction as the tangent vector to
$\vec{r}(t)$.
}

\content{
\begin{definition}
Given a vector valued function $\vec{r}(t)$, the unit normal vector is given by 
$$ \vec{N}(t) = \frac{\vec{T}'(t)}{|\vec{T}'(t)|}. $$
\end{definition}
}{% presentation
\vspace{2in}
}{% nopresentation
\vspace{1in}
}{% comments
Since $|\vec{T}(t)|=1$ for all $t$, $\vec{T}'(t)$ is orthogonal to $\vec{T}(t)$.
Thus this vector can be thought of as the direction that the particle is turning
at time $t$.
}

\content{
\begin{definition}
Given a vector valued function $\vec{r}(t)$, the binormal vector $\vec{B}(t)$ is
given by 
$$ \vec{B}(t) = \vec{T}(t)\times \vec{N}(t). $$
\end{definition}
}{% presentation
\vspace{2in}
}{% nopresentation
\vspace{1in}
}{% comments
This vector is orthogonal to both $\vec{T}(t)$ and $\vec{N}(t)$, and is also a
unit vector. 

Draw a picture of how this gives a basis (called the TNB basis), which is
important to differential geometry and applications to the motion of spacecraft. 
}

\content{
\begin{example} Find the unit normal vector and binormal vectors for the
function 
$$ \vec{r}(t) = \la \cos t, \sin t, t\ra. $$
\end{example}
}{% presentation
\vspace{5in}
}{% nopresentation
\vspace{3in}
}{% comments
}

\subsection{Motion in Space}

\content{

}{% presentation
\vspace{3in}
}{% nopresentation
\vspace{2in}
}{% comments
demonstrate how the derivative of a vector valued function is the tanged vector,
which is the direction of motion of a particle traveling along $\vec{r}(t)$. 

Similarly, the second derivative is the acceleration, just as in the calc I
case. 

Mention that the speed of the particle at a point is $|\vec{r}'(t)|$.
}

\content{
\begin{example}
A moving particle starts at an initial position of $\vec{r}(0) = \la 1,0,0\ra$,
and with initial velocity $\vec{v}(0) = \la 1,-1,1\ra$.  It's acceleration is
given by $\vec{a}(t) = \la 4t, 6t, 1 \ra$. Find its velocity and position
functions. 
\end{example}
}{% presentation
\vspace{4in}
}{% nopresentation
}{% comments
Mention that there are good applications of vector-valued functions to
projectile motion and Kepler's Laws of Planetary Motion in the book if you are
interested in reading them. 
}

